% Vietnamese translation for OI Public Library
% Source-PDF: IOI中国国家候选队论文/2008/Day1/10.俞华程《矩阵乘法在信息学中的应用》/slide.pdf
\documentclass[11pt]{article}
\usepackage{tikz}
\usetikzlibrary{arrows.meta}
\usepackage{oipl-vn}
\renewcommand{\figurename}{Hình}

\title{Ứng dụng của phép nhân ma trận trong tin học}
\author{Du Hoa Trình}
\date{}

\begin{document}
\maketitle
\sourcepath{IOI 2008 / Yu Huacheng / Matrix multiplication applications}

\begin{abstract}
Bản gốc là bộ slide đi kèm bài viết chính về ứng dụng của phép nhân ma trận.
Bản tiếng Việt này chuyển nội dung slide thành một bài ghi chú có cấu trúc:
giữ các lập luận, công thức, ví dụ thuật toán cốt lõi và các hình trực quan
chính của phần trình chiếu.
\end{abstract}

\section{Các hướng ứng dụng}

Phép nhân ma trận thường xuất hiện khi cần:
\begin{itemize}
  \item tối ưu quy hoạch động hoặc tăng tốc mô phỏng;
  \item đếm đường đi thông qua ma trận kề của đồ thị;
  \item kết hợp với chia đôi hoặc chia để trị.
\end{itemize}

Tài liệu này tập trung vào cách dùng ma trận kề và lũy thừa ma trận để
biểu diễn số đường đi trong đồ thị, rồi áp dụng vào một bài toán có
trạng thái biến đổi theo thời gian.

\section{Lũy thừa của ma trận kề}

Xét đồ thị có $N$ đỉnh và ma trận kề $G$. Với đồ thị không trọng số,
$G[a,b]=1$ khi có cạnh từ $a$ đến $b$, ngược lại bằng $0$. Theo định
nghĩa phép nhân ma trận,
\[
  G^2[a,b]=\sum_{i=1}^{N}G[a,i]G[i,b].
\]
Tích $G[a,i]G[i,b]$ bằng $1$ khi và chỉ khi có đường đi
$a\to i\to b$. Vì vậy $G^2[a,b]$ chính là số đường đi độ dài $2$ từ
$a$ đến $b$.

Tương tự,
\[
\begin{aligned}
G^3[a,b]
  &=\sum_{i=1}^{N}G[a,i]G^2[i,b] \\
  &=\sum_{i=1}^{N}G[a,i]\left(\sum_{j=1}^{N}G[i,j]G[j,b]\right) \\
  &=\sum_{i=1}^{N}\sum_{j=1}^{N}G[a,i]G[i,j]G[j,b].
\end{aligned}
\]
Mỗi tích khác $0$ tương ứng với một đường đi
$a\to i\to j\to b$, nên $G^3[a,b]$ là số đường đi độ dài $3$.

\begin{figure}[ht]
\centering
\begin{tikzpicture}[
  >=Stealth,
  every node/.style={font=\small},
  v/.style={circle,draw,minimum size=7mm,inner sep=1pt},
  arc/.style={->,thick}
]
\node[v] (a) at (0,0) {$a$};
\node[v] (i) at (1.4,0) {$i$};
\node[v] (b) at (2.8,0) {$b$};
\draw[arc] (a) -- (i);
\draw[arc] (i) -- (b);
\node at (1.4,-.65) {độ dài \(2\)};
\node[v] (a2) at (4.2,0) {$a$};
\node[v] (i2) at (5.4,0) {$i$};
\node[v] (j2) at (6.6,0) {$j$};
\node[v] (b2) at (7.8,0) {$b$};
\draw[arc] (a2) -- (i2);
\draw[arc] (i2) -- (j2);
\draw[arc] (j2) -- (b2);
\node at (6,-.65) {độ dài \(3\)};
\end{tikzpicture}
\caption{Hai hình minh họa trong slide: tích trong \(G^2[a,b]\) đếm đường \(a\to i\to b\), còn tích trong \(G^3[a,b]\) đếm đường \(a\to i\to j\to b\).}
\end{figure}

Suy rộng, $G^k[a,b]$ là số đường đi độ dài $k$ từ $a$ đến $b$. Kết luận
này không cần xử lý riêng cạnh song song hay vòng tự khép; chúng đều
được tính đúng trong ma trận kề tương ứng.

\section{Bài toán mô phỏng trên đồ thị}

Cho một đồ thị vô hướng có không quá $50$ đỉnh, một đỉnh bắt đầu $s$ và
một đỉnh kết thúc $t$. Mỗi đơn vị thời gian người chơi di chuyển một
lần theo cạnh. Một số kẻ săn mồi di chuyển tuần hoàn, chu kỳ không quá
$4$. Người chơi không được đi vào vị trí có kẻ săn mồi. Cần tính số
cách đi từ $s$ đến $t$ tại thời điểm $u$, với $u\le 2\cdot 10^9$.

\begin{figure}[ht]
\centering
\begin{tikzpicture}[
  every node/.style={font=\small},
  v/.style={circle,fill=black,inner sep=2.1pt},
  safe/.style={line width=.5pt},
  danger/.style={line width=1pt,densely dashed}
]
\node[v,label=above left:{\(8\)}] (a) at (0,1.1) {};
\node[v,label=above:{\(s\)}] (s) at (1.55,1.1) {};
\node[v] (r) at (2.7,1.45) {};
\node[v] (m) at (1.05,-.35) {};
\node[v,label=right:{\(t\)}] (t) at (2.4,-.15) {};
\node[v,red] (p0) at (-.7,-.2) {};
\draw[safe] (a) -- (s) -- (r) -- (m) -- (a);
\draw[safe] (s) -- (m);
\draw[safe] (s) -- (t);
\draw[danger,red] (p0) -- (a);
\draw[danger,blue] (m) -- (r);
\draw[danger,blue] (m) -- (t);
\draw[danger,blue] (t) -- (r);
\node at (1.1,-1.05) {\(u=4\), thời điểm \(0\)};
\end{tikzpicture}
\caption{Ảnh chụp trực quan của bài \textit{Swamp Crocodile}: tại mỗi thời điểm, các đỉnh/cạnh bị kẻ săn mồi chiếm làm thay đổi ma trận chuyển \(G_u\).}
\end{figure}

Nếu không có kẻ săn mồi, lời giải trực tiếp là dùng $G^u$: số cách từ
$s$ đến $t$ là $G^u[s,t]$. Khi có kẻ săn mồi, ma trận chuyển trạng thái
phụ thuộc vào thời gian, vì ở mỗi thời điểm một số đỉnh bị cấm.

Gọi $A_u[a,b]$ là số cách để ở đỉnh $a$ tại thời điểm $0$ và ở đỉnh
$b$ tại thời điểm $u$. Nếu không có cấm đỉnh thì $A_u=G^u$. Khi có cấm
đỉnh, đặt $G_u$ là ma trận chuyển trạng thái tại bước thời gian $u$:
các cạnh đi vào đỉnh bị kẻ săn mồi chiếm tại thời điểm đó được xóa khỏi
ma trận. Khi đó
\[
  A_u=A_{u-1}G_u=G_1G_2\cdots G_u.
\]

\section{Khai thác tính chu kỳ}

Chu kỳ của mỗi kẻ săn mồi không quá $4$, nên trạng thái cấm đỉnh lặp
lại với chu kỳ
\[
  \operatorname{lcm}(1,2,3,4)=12.
\]
Do đó $G_i=G_{i+12}$. Đặt
\[
  G'=G_1G_2\cdots G_{12},\qquad u=12p+q,\qquad 0\le q<12.
\]
Ta có
\[
  G_1G_2\cdots G_u=(G')^pG_1G_2\cdots G_q.
\]
Vì vậy chỉ cần tính trước tích một chu kỳ $G'$ và tích tiền tố
$G_1\cdots G_q$, rồi dùng lũy thừa nhanh cho $(G')^p$.

\section{Độ phức tạp}

Với phép nhân ma trận thường, nhân hai ma trận $N\times N$ mất
$O(N^3)$. Các bước của thuật toán gồm:
\begin{itemize}
  \item tính $G_1G_2\cdots G_{12}$ trong $O(N^3)$ vì $12$ là hằng số;
  \item tính tiền tố $G_1\cdots G_q$ trong $O(N^3)$;
  \item lũy thừa nhanh $(G')^p$ trong $O(N^3\log p)=O(N^3\log u)$;
  \item nhân các ma trận cuối để lấy đáp án trong $O(N^3)$.
\end{itemize}
Tổng độ phức tạp là $O(N^3\log u)$, phù hợp với $N\le 50$ và
$u\le 2\cdot 10^9$.

\section{Tổng kết}

Phép nhân ma trận xử lý được nhiều dạng bài toán đếm đường đi:
\begin{itemize}
  \item số đường đi có số cạnh cố định giữa hai đỉnh;
  \item số đường đi có số cạnh nằm trong một khoảng;
  \item số đường đi có số cạnh cố định giữa mọi cặp đỉnh;
  \item số đường đi có số cạnh nằm trong một khoảng giữa mọi cặp đỉnh.
\end{itemize}

Ý tưởng chung là biểu diễn một bước chuyển bằng ma trận, sau đó dùng
phép nhân và lũy thừa nhanh để ghép nhiều bước chuyển. Khi ma trận thay
đổi tuần hoàn theo thời gian, ta nhóm cả chu kỳ thành một khối chuyển
trạng thái rồi tiếp tục dùng lũy thừa ma trận.

\end{document}
